\documentclass{article}
\usepackage[utf8]{inputenc}
\usepackage{geometry}
\geometry{a4paper, margin=1in}

\title{What if Forgiveness is an Error-Correction Protocol?}
\author{Albert Jan van Hoek}
\date{\today}

\begin{document}

\maketitle

\section*{Theory of Long-term Collaboration (TLC)}

I’ve been developing what I’ll call the \textbf{Theory of Long-term Collaboration (TLC)} — a way to understand how intelligences work together over time.

Humans are intelligent. Brains are networks. Intelligence is the capacity of a network to keep updating itself — a recursive learning loop.

If each intelligence is a learning network, collaboration isn’t just connecting two dots. It’s connecting two updating systems.

That means the link between intelligences — the edge — isn’t a thin line. It’s a shared space where alignment, coordination, repair, and learning happen together: a small meta-network co-created by the participants.

Importantly, this meta-network is itself a learning space. Over time, not only what is done changes, but how both work together changes too.

\section*{Patterns of Behavior}

Viewed this way, patterns of behavior start to explain themselves. Some behaviors optimize this meta-network. Others create friction.

Forgiveness, for instance, efficiently restores collaboration after conflict.

From this perspective, norms don’t need to be imposed from the outside. They emerge from what sustains collaboration over time.

\section*{Two Kinds of Freedom}

This also suggests two kinds of freedom:
\begin{itemize}
    \item The familiar freedom: the ability to think, speak, and act as you wish.
    \item A systemic freedom: the freedom to act within the constraints of sustaining the network. Some choices simply make sustained collaboration harder.
\end{itemize}

TLC doesn’t replace ethics; it reframes them. Ethics emerge, repeatedly, within any collaboration between autonomous intelligent agents. TLC offers a neutral, precise, and technical way to talk about collaboration — turning “be a good collaborator” into something you can analyze, measure, and improve.

\section*{Maintenance Operations}

Forgiveness becomes an error-correction protocol. \\
Honesty becomes low-latency signal transmission. \\
Gratitude becomes reinforcement that strengthens the meta-network.

These aren’t just virtues — they’re maintenance operations.

\section*{Relevance for Public Health}

Details may vary by substrate — human, AI, or other intelligent agents — but TLC offers a new lens to understand collaboration.

I think TLC is relevant for public health. Society can be seen as a mega-meta-network, and therefore public health isn't only about bodies — it's about the quality of the connections between them. Loneliness, conflict, and broken trust aren’t just social problems; they're network degradation. And interventions that restore collaboration — at any scale — are, in this framework, a form of medicine.

\end{document}
